% Modelo de prova da Maratona de Programação InterFatecs
% Autor: Leandro Luque
% Versão: 1.0
\documentclass[11pt]{problemset}
\usepackage[portuges]{babel}
\usepackage{graphicx,subfigure}

%%%%%%%%%%%%%%%%%%%%%%%%%%%%%%%%%%%
\usepackage{tikz} % Diagramas.
\tikzset{
  treenode/.style = {shape=rectangle, rounded corners,
                     draw, align=center,
                     top color=white, bottom color=blue!20},
  root/.style       = {treenode, font=\ttfamily\normalsize},
  opcao/.style      = {treenode, font=\ttfamily\normalsize},
  decisao/.style    = {treenode, draw}
}

\contestname{Maratona de Programação InterFatecs}
\contestshortname{Maratona de Programação InterFatecs 2021 - }
% Descomente para inserir o logo da InterFatecs nas páginas dos problemas
\headerlogo{logo}

\setlength{\parindent}{0.0in} 
\setlength{\parskip}{0.1in}

\begin{document}
% Capa
\begin{figure}[h!]
	\centering
    \includegraphics[width=1\textwidth]{CapaFinal.png}
\end{figure}

%\footnotesize

\section{Instruções}

Este caderno contém 11 problemas – identificados por letras de A até K, com páginas numeradas de 3 até \pageref{LastPage}. Verifique se seu caderno está completo.

Informações gerais

\begin{enumerate}
\item Sobre a competição
\begin{enumerate}
\item A competição possui duração de 5 horas (início as 13:00h término as 18:00h);
%\item NÃO é permitido acesso a conteúdo da Internet ou qualquer outro meio eletrônico digital;
%\item É permitido somente acesso a conteúdo impresso (cadernos, apostilas, livros);
\item É permitido a consulta a materiais já publicados anteriormente ao dia da competição
\item Não é permitido a comunicação com o técnico ou qualquer outra pessoa que não seja a equipe para tirar dúvidas sobre a maratona
\item É vedada a comunicação entre as equipes durante a competição, bem como a troca de material de consulta entre elas;
%\item Cada equipe terá acesso a 1 computador dotado do ambiente de submissão de programas (BOCA), dos compiladores, link-editores e IDEs requeridos pelas linguagens de programação permitidas;
\item Cada integrante da equipe poderá utilizar o seu computador/notebook para resolver os problemas, porém, apenas um competidor da equipe ficará responsável pela submissão
%\item NÃO é permitido o uso de notebooks ou outro tipo de computador ou assistente pessoal;
\item Os problemas têm o mesmo valor na correção. 
\end{enumerate}
\item Sobre o arquivo de solução e submissão:
  \begin{enumerate}
  \item O arquivo de solução (o programa fonte) deve ter o mesmo nome que o especificado no enunciado (logo após o título do problema);
  \item confirme se você escolheu a linguagem correta e está com o nome de arquivo correto antes de submeter a sua solução;
  \item NÃO insira acentos no arquivo-fonte.
  \end{enumerate}
\item Sobre a entrada
  \begin{enumerate}
  \item A entrada de seu programa deve ser lida da entrada padrão (não use interface gráfica);
  \item Seu programa será testado em vários casos de teste válidos além daqueles apresentados nos exemplos. Considere que seu programa será executado uma vez para cada caso de teste.
  \end{enumerate}
\item Sobre a saída
  \begin{enumerate}
  \item A saída do seu programa deve ser escrita na saída padrão;
  \item Não exiba qualquer outra mensagem além do especificado no enunciado.
  \end{enumerate}
\end{enumerate}
% Inclui os problemas na prova.

% Cada problema deve ser inserido em uma pasta diferente.

% Problemas da 1ª Fase
%\enincludeproblem{parking}
%\includeproblem{xuquisim}
%\includeproblem{epidemia}
%\includeproblem{taylor}
%\includeproblem{celulas}
%\includeproblem{lojinha}
%\includeproblem{digitos}
%\includeproblem{nakitomon}
%\includeproblem{megabobagem}
%\includeproblem{jogovida}
%\includeproblem{bochas}

% Problemas da Fase Final
\enincludeproblem{abbreviating}
\includeproblem{agrupamentos}
%\includeproblem{sertanejo}
%\includeproblem{chinelandia}
%\includeproblem{fifo}
%\includeproblem{escola}
%\includeproblem{festa}
%\includeproblem{quociente}
%\includeproblem{reforma}
%\includeproblem{acessibilidade}
%\includeproblem{lubank}
%\includeproblem{orquestrador}
\end{document}
